% Preparación autónoma de VII, 2026-09-30. Fragmento sin inputs adicionales.
% RESULTADOS_RECUPERADOS, pruebas reunidas sin modificación de los originales:
% /Users/ruben/Documents/excelencia academica/output/CUMPLIMIENTO_EDITORIAL_20260928/fuentes/VIII_ES/manuscrito/29_retorno_y_memoria_traslacional.tex
% /Users/ruben/Documents/excelencia academica/output/CUMPLIMIENTO_EDITORIAL_20260928/fuentes/VIII_ES/manuscrito/30_pantalla_y_respuesta.tex
% /Users/ruben/Documents/excelencia academica/output/CUMPLIMIENTO_EDITORIAL_20260928/fuentes/VIII_ES/manuscrito/30b_variacion_energia_memoria.tex
% /Users/ruben/Documents/excelencia academica/output/CUMPLIMIENTO_EDITORIAL_20260928/fuentes/VIII_ES/manuscrito/30c_composicion_corriente_conexion.tex
% /Users/ruben/Documents/excelencia academica/output/CUMPLIMIENTO_EDITORIAL_20260928/fuentes/VIII_ES/manuscrito/30d_realizacion_geometrica.tex
% /Users/ruben/Documents/excelencia academica/output/CUMPLIMIENTO_EDITORIAL_20260928/fuentes/VIII_ES/manuscrito/31_corriente_y_respuesta_cartan.tex
% /Users/ruben/Documents/excelencia academica/output/CUMPLIMIENTO_EDITORIAL_20260928/fuentes/VIII_ES/manuscrito/33_corriente_espinorial_y_densidad.tex
% Dependencias internas ya presentes: núcleo común, rutas, lectores y unidades.
% Las inversas de 31 están demostradas en 10, hmtcuatro:gea:torsion.
% Legendre, fuente completa y restricciones permanecen en 20.
% Se incluye de 33 sólo lo utilizado por 10+20, no sus estados cosmológicos.
\section{Antecedentes geométricos y materiales de la acción conjunta}
\label{hmtcuatro:ant:seccion}

Las semillas y las operaciones de APP, la orientación y el régimen
del TRIT y el transporte del TPK producen las rutas enriquecidas sobre
las que actúan los lectores del núcleo común. Se conservan conjuntamente
hoja, orientación, acarreo, frontera y memoria; el retorno de fase no
reinicia ese estado. Las construcciones del continuo pertenecen al mismo
objeto y no se convierten aquí en sectores independientes. La
realización que sigue transporta a la geometría y a la materia esa
estructura ya construida. Sus unidades y acoplamientos son salidas
HMT anteriores: ninguna comparación numérica interviene como generador.

Para que la acción y su reducción puedan leerse dentro de este
manuscrito, reunimos las operaciones de pantalla, la variación efectiva
del transporte y la corriente espinorial que aquéllas utilizan. La
eliminación de Cartan--Holst se demuestra en
\ref{hmtcuatro:gea:torsion}; la acción total y su reducción canónica,
en \ref{hmtcuatro:accion-retroaccion-restricciones}. Las fórmulas de
esta sección especifican los datos y los cálculos que entran en ambas.

\subsection{Incidencia, pantalla y soldadura}
\label{hmtcuatro:ant:pantalla}

La representación cúbica de la frontera de una ruta tiene las seis
caras \(F=\{(i,\epsilon):1\leq i\leq3,\ \epsilon=\pm1\}\).
Sobre \(H_F=\mathbb R^F\), la oposición es
\(Re_{i,\epsilon}=e_{i,-\epsilon}\). Con
\[
 u=\sum_{i,\epsilon}e_{i,\epsilon},\qquad
 P_u=\frac{uu^{\mathsf T}}6,\qquad
 P_-=\frac{I-R}2,\qquad
 P_0=\frac{I+R}2-P_u ,
\]
los tres proyectores son ortogonales, suman \(I\) y tienen rangos
\(1,3,2\). La incidencia entre caras y vértices
\(\nu\in\{-1,1\}^3\) es
\(M_{(i,\epsilon),\nu}=\mathbf1_{\{\nu_i=\epsilon\}}\).

\begin{proposition}[Pantalla de la incidencia orientada]
\label{hmtcuatro:ant:prop-pantalla}
Se cumple
\[
 MM^{\mathsf T}=12P_u+4P_-,\qquad
 \ker M^{\mathsf T}=P_0H_F .
\]
El cociente \(Q=H_F/P_0H_F\), identificado con
\((P_u+P_-)H_F\), es de dimensión cuatro. La forma
\(B=P_--P_u\) sobre \(Q\) tiene firma \((-+++)\).
\end{proposition}
\begin{proof}
Cada cara contiene cuatro vértices; dos caras opuestas no comparten
ninguno y dos caras de ejes distintos comparten dos. Por tanto,
\[
 MM^{\mathsf T}=4I+2(uu^{\mathsf T}-I-R)
              =12P_u+4P_- .
\]
La identidad
\(\|M^{\mathsf T}x\|^2=\langle x,MM^{\mathsf T}x\rangle\)
determina el núcleo. En la base ortonormal
\[
 t=u/\sqrt6,\qquad d_i=(e_{i,+}-e_{i,-})/\sqrt2
\]
de \(Q\), la matriz de \(B\) es \(\operatorname{diag}(-1,1,1,1)\).
\end{proof}

Sea \(W_\square=\sqrt6P_u+\sqrt2P_-\). Entonces
\[
 W_\square e_{i,\epsilon}=t+\epsilon d_i=:n_{i,\epsilon},
 \qquad B(n_{i,\epsilon},n_{i,\epsilon})=0,\qquad
 MM^{\mathsf T}|_Q=2W_\square^*W_\square|_Q .
\]
Estas igualdades se obtienen aplicando los proyectores a la base de
caras. Las rotaciones propias del cubo preservan los proyectores,
la orientación temporal \(t\) y la forma \(B\).

La ruta enriquecida proporciona a cada arista \(a\) su cara,
su signo \(\sigma_m(a)\) y el transporte
\(U_m(a):Q_{s(a)}\to Q_{t(a)}\). Su soldadura celular es
\begin{equation}
 \beta_m(a)=\ell_m\sigma_m(a)n_{\lambda(\underline a)}^{s(a)},
 \quad \ell_m=\ell_0\,9^{-m},\quad
 \beta_m(\bar a)=-U_m(a)\beta_m(a).
 \label{hmtcuatro:ant:soldadura}
\end{equation}
La inversión \(U_m(\bar a)=U_m(a)^{-1}\) implica que invertir dos
veces recupera la soldadura. Se conserva además el registro ordenado
del acarreo: no se sustituye por el vector de llegada.

\begin{proposition}[Naturalidad de la soldadura]
\label{hmtcuatro:ant:refinamiento-soldadura}
Sea \(a=a_1\cdots a_9\) un refinamiento compatible: las incidencias
hijas transportadas al origen del predecesor conservan cara y signo,
y \(\ell_{m+1}=\ell_m/9\). Si
\(\mathcal U_j:Q_{m,s(a)}\to Q_{m+1,s(a_j)}\) incluye la
identificación por refinamiento y el transporte hasta \(s(a_j)\),
entonces
\[
 \beta_m(a)=\sum_{j=1}^9\mathcal U_j^{-1}\beta_{m+1}(a_j).
\]
\end{proposition}
\begin{proof}
Cada sumando en la fibra de origen es
\((\ell_m/9)\sigma_m(a)n_{\lambda(\underline a)}\). Su suma da
\eqref{hmtcuatro:ant:soldadura}. Invertir la ruta invierte el orden
de los transportes y cada transporte; la fórmula de inversión de
la soldadura da el signo global. La concatenación conserva en el
mismo orden la memoria de las nueve incidencias. La igualdad se
itera por composición de refinamientos compatibles.
\end{proof}

\subsection{Representación del transporte y realización geométrica}
\label{hmtcuatro:ant:geometria}

El registro de una ruta \(\gamma\) cuenta las inversiones y los
cambios efectivos de hoja mediante el cociclo
\(c(\gamma)=(c_g(\gamma),c_s(\gamma))\in\mathbb Z^2\).
El recuento es acumulativo en las rutas hacia adelante; su
extensión al grupoide es orientada y asigna incrementos opuestos
a la ruta inversa.
Para rutas composables,
\(c(\gamma_2\gamma_1)=c(\gamma_2)+c(\gamma_1)\), y
\(c(\bar\gamma)=-c(\gamma)\). Sea \(R_4\) un elemento de orden
cuatro en \(\operatorname{Spin}^+(3,1)\) cuya representación
vectorial \(\Lambda(R_4)\) fija el vector temporal unitario \(u_0\).
La representación afín recibida es
\[
 \rho_{\rm C}^{\rm disc}(\gamma)
 =\bigl(R_4^{c_g(\gamma)},\ell_0c_s(\gamma)u_0\bigr).
\]
Con el producto
\((R,a)(S,b)=(RS,a+\Lambda(R)b)\), la fijación de \(u_0\) y la
aditividad de \(c\) prueban directamente que \(\rho_{\rm C}^{\rm disc}\)
conserva identidad, concatenación e inversión. Cuando la componente
rotacional retorna a la identidad, la traslación de memoria permanece
determinada por \(c_s(\gamma)\), que no se reinicia.

La normalización
\[
 N_{\ell_0}(R,a)=(\Lambda(R),a/\ell_0),\qquad
 \rho_{\rm C}=N_{\ell_0}\circ\rho_{\rm C}^{\rm disc}
\]
también es un homomorfismo, pues
\((a+\Lambda(R)b)/\ell_0=a/\ell_0+\Lambda(R)(b/\ell_0)\).
El levantamiento \(R_4^{c_g(\gamma)}\) se conserva cuando actúan
espinores; no se recupera sólo a partir de \(\Lambda(R_4^{c_g(\gamma)})\),
cuya representación tiene núcleo \(\{1,-1\}\).

En una realización de pantalla \(J_{{\rm scr},m}\) que transporta
la forma y la soldadura, \(e_m=J_{{\rm scr},m}\beta_m\).
La realización suave considerada consiste en una región orientada
\(\mathcal U\) de dimensión cuatro, una cotetrada no degenerada
\(e\), una conexión de Lorentz \(\omega\) y una realización de rutas
\(\mathfrak R_{\rm C}\) con la condición de compatibilidad
\begin{equation}
 \operatorname{Hol}_{\mathcal A_{\ell_0}}
          (\mathfrak R_{\rm C}(\gamma))
 =\rho_{\rm C}(\operatorname{Hol}_{\rm TPK}(\gamma)),\qquad
 \mathcal A_{\ell_0}=
 \begin{pmatrix}\omega&\ell_0^{-1}e\\0&0\end{pmatrix}.
\label{hmtcuatro:ant:compatibilidad-holonomia}
\end{equation}
El argumento de \(\rho_{\rm C}\) conserva el cociclo de la
holonomía enriquecida; no se sustituye por su sola fase reducida.
La escala \(\ell_0>0\) es constante en la coordenatización. Éstos
son los datos y la condición de la realización suave utilizada;
preservar la composición de caminos no identifica entre sí caminos
distintos de iguales extremos ni elimina su memoria. La métrica
es \(g=\eta_{IJ}e^I\otimes e^J\), \(\eta=\operatorname{diag}(-1,1,1,1)\).

\begin{proposition}[Curvatura, covariancia y Bianchi]
\label{hmtcuatro:ant:curvatura}
Con \(F_\omega=\mathrm d\omega+\omega\wedge\omega\) y
\(T=\mathrm de+\omega\wedge e\),
\[
 \mathcal F_{\ell_0}=
 \begin{pmatrix}F_\omega&\ell_0^{-1}T\\0&0\end{pmatrix},
 \qquad D_\omega T=F_\omega\wedge e,\quad D_\omega F_\omega=0 .
\]
Un cambio de marco \(g_L:\mathcal U\to SO^+(3,1)\) actúa por
\(e'=g_Le\),
\(\omega'=g_L\omega g_L^{-1}-\mathrm dg_L\,g_L^{-1}\);
entonces \(T'=g_LT\), \(F_{\omega'}=g_LF_\omega g_L^{-1}\).
\end{proposition}
\begin{proof}
La multiplicación de matrices de formas da el bloque diagonal
\(\mathrm d\omega+\omega\wedge\omega\) y el traslacional
\(\ell_0^{-1}(\mathrm de+\omega\wedge e)\). Al sustituir los
campos transformados, los términos con \(\mathrm dg_L\) se
cancelan. Expandir \(D_\omega T\) deja
\((\mathrm d\omega+\omega\wedge\omega)\wedge e\).
En \(D_\omega F_\omega\), las derivadas de \(\omega\) y los
productos cúbicos se cancelan por pares.
\end{proof}

La hoja trítica \(\tau\in\{+1,0,-1\}\) conserva su régimen en la
realización de Cartan con generadores \(J_{ab},P_a^{(\tau)}\),
acción vectorial de Lorentz y
\([P_a^{(\tau)},P_b^{(\tau)}]=-\tau J_{ab}\). Para
\(\mathcal A_\tau=\frac12\omega^{ab}J_{ab}+\ell^{-1}e^aP_a^{(\tau)}\),
\[
 \mathcal F_\tau=
 \tfrac12(F_\omega^{ab}-\tau\ell^{-2}e^a\wedge e^b)J_{ab}
 +\ell^{-1}T^aP_a^{(\tau)} .
\]
En efecto, el corchete de Lorentz da \(F_\omega\), el mixto
da \(D_\omega e\) y el traslacional da el término
\(-\tau e\wedge e/(2\ell^2)\). La condición
\eqref{hmtcuatro:ant:compatibilidad-holonomia} corresponde a la
hoja central; las hojas extremas utilizan su representación
\(\rho_{{\rm C},\tau}\) y su condición de holonomía propias.

\subsection{Diferencial del transporte y corriente material}
\label{hmtcuatro:ant:corriente-transporte}

En una realización celular finita, sea
\(U_a:\mathcal H_{s(a)}\to\mathcal H_{t(a)}\) un transporte
invertible entre fibras hermíticas. No se presupone que todo
transporte de Lorentz sea unitario para el producto energético.
Se toma una orientación por arista. Para
\[
 (Df)_a=f_{t(a)}-U_af_{s(a)},\quad r=Df,\quad
 q=\mathsf W r,\quad v_a=U_af_{s(a)},\quad
 E=\langle r,\mathsf W r\rangle ,
\]
el peso \(\mathsf W=\mathsf W^*>0\) puede acoplar aristas.
Este peso no es \(W_\square\). Las variaciones corresponden a
la representación de la conexión y al problema de borde elegido.

\begin{proposition}[Variación completa]
\label{hmtcuatro:ant:variacion-completa}
Para una curva diferenciable de datos, con
\(A_a=\dot U_aU_a^{-1}\), se cumple
\begin{equation}
 \dot E=2\operatorname{Re}\sum_a
 \langle q_a,\dot f_{t(a)}-U_a\dot f_{s(a)}-A_av_a\rangle
 +\langle r,\dot{\mathsf W}r\rangle .
 \label{hmtcuatro:ant:variacion}
\end{equation}
A campo y peso fijos, el gradiente completo del transporte es
\(G_a=-2q_av_a^*\):
\[
 \delta_U E=\sum_a\operatorname{Re}\operatorname{Tr}(G_a^*A_a).
\]
En la subclase unitaria \(A_a^*=-A_a\), basta su parte
antihermítica \(J_a=v_aq_a^*-q_av_a^*\).
\end{proposition}
\begin{proof}
La regla del producto da
\(\dot E=2\operatorname{Re}\langle\mathsf W r,\dot r\rangle+
\langle r,\dot{\mathsf W}r\rangle\).
Derivar cada residuo produce \eqref{hmtcuatro:ant:variacion}.
Como \(\operatorname{Tr}(v_aq_a^*A_a)=q_a^*A_av_a\), se obtiene
\(G_a\). La parte hermítica de \(G_a\) es ortogonal, para
\(\operatorname{Re}\operatorname{Tr}(X^*Y)\), a las variaciones
antihermíticas. Esta reducción sólo se aplica en esa subclase.
\end{proof}

Para una ruta con \(U_\gamma=U_N\cdots U_1\), sea
\(B_k=U_N\cdots U_{k+1}\). La regla del producto demuestra
\begin{equation}
 \dot U_\gamma U_\gamma^{-1}
 =\sum_k B_kA_kB_k^{-1}.
 \label{hmtcuatro:ant:variacion-ruta}
\end{equation}
Consecuentemente un gradiente \(G_\gamma\) se transporta hacia
la arista \(k\) como
\(G_k=B_k^*G_\gamma(B_k^{-1})^*\). La ciclicidad de la traza
prueba esta fórmula incluso si \(B_k\) no es unitario.

Si las coordenadas reales de variación de conexión son
\(\theta_b\), se escribe su diferencial efectivo como
\(A_a(\theta)=\sum_b B_{ab}\theta_b\). Para
\(S_{\rm mem}=\mathfrak a E\), a campo, peso y sección de acción
\(\mathfrak a\) fijos,
\begin{equation}
 \delta S_{\rm mem}=-\tfrac12\sum_b\theta_b s_b,\qquad
 s_b=4\mathfrak a\,\operatorname{Re}
          \sum_a\langle q_a,B_{ab}v_a\rangle .
 \label{hmtcuatro:ant:corriente-celular}
\end{equation}
Sustituir \(A_a(\theta)\) en la variación demuestra la igualdad
y fija el factor \(4\). Si el emparejamiento celular orientado es
una matriz invertible \(\mathsf M\), definida por
\(\delta S_{\rm mem}=-\frac12\theta^{\mathsf T}\mathsf M\sigma\),
la corriente es \(\sigma=\mathsf M^{-1}s\). Esta matriz conserva
los volúmenes y la orientación; no puede sustituirse por una
división escalar sin fijar una base en la que corresponda.
Si también varían \(f,\mathsf W,\mathfrak a\), se conservan todos
los términos de \eqref{hmtcuatro:ant:variacion} y \(E\delta\mathfrak a\).

Para un peso por arista, su inversión conserva la respuesta completa
al transformar conjuntamente
\[
 U_{\bar a}=U_a^{-1},\qquad
 r_{\bar a}=-U_a^{-1}r_a,\qquad
 \mathsf W_{\bar a}=U_a^*\mathsf W_aU_a .
\]
La sustitución da
\(r_{\bar a}^*\mathsf W_{\bar a}r_{\bar a}=r_a^*\mathsf W_ar_a\).
Su derivada conserva necesariamente
\(\dot{\mathsf W}_{\bar a}
=U_a^*(A_a^*\mathsf W_a+\dot{\mathsf W}_a+\mathsf W_aA_a)U_a\):
eliminar este término alteraría la corriente. Si el peso acopla varias
aristas, se aplica la misma congruencia a la matriz completa con el
transporte diagonal por bloques; se conservan también sus bloques
cruzados.

\begin{proposition}[Compatibilidad variacional por refinamiento]
\label{hmtcuatro:ant:refinamiento-corriente}
Para mapas fijos \(J_m,K_m\), supóngase que, a lo largo de la
variación considerada,
\[
 D_{m+1}J_m=K_mD_m,\qquad
 K_m^*\mathsf W_{m+1}K_m=\mathsf W_m .
\]
Entonces \(E_{m+1}(J_mf)=E_m(f)\) y sus diferenciales coinciden.
Si además las acciones coinciden y \(P_m\) transporta las
variaciones de conexión al refinamiento, sus corrientes cumplen
\[
 s_m=P_m^{\mathsf T}s_{m+1},\qquad
 \mathsf M_m\sigma_m=P_m^{\mathsf T}\mathsf M_{m+1}\sigma_{m+1}.
\]
\end{proposition}
\begin{proof}
Sustituir \(D_{m+1}J_m=K_mD_m\) en la energía y después la
congruencia del peso prueba la primera igualdad para toda la
curva. Derivarla y utilizar
\(-\frac12\theta^{\mathsf T}s_m
=-\frac12(P_m\theta)^{\mathsf T}s_{m+1}\)
prueba las restantes. Si \(J_m\) varía, su diferencial aporta
\(2\operatorname{Re}\langle D_{m+1}J_mf,
\mathsf W_{m+1}D_{m+1}\dot J_m f\rangle\); no se descarta.
\end{proof}

La corriente de memoria así obtenida procede de su acción.
La corriente espinorial siguiente procede de la acción afín de
los campos. Se componen cuando son términos de una misma acción;
no se identifican sólo por compartir el transporte. En particular,
eliminar una contorsión de la que dependen también el campo
preparado, la incidencia o su mínimo exige conservar esos
diferenciales. La fórmula lineal de Cartan utilizada a continuación
corresponde al problema material afín expresamente declarado.

\subsection{Representación espinorial y corriente afín}
\label{hmtcuatro:ant:espinores}

Se fija la estructura de espín de la realización geométrica. Las
matrices heredadas, después de complejificar el módulo real, son
\begin{equation}
 \begin{gathered}
 J=\begin{pmatrix}0&-1\\1&0\end{pmatrix},\quad
 R=\begin{pmatrix}0&1\\1&0\end{pmatrix},\quad
 Z=\begin{pmatrix}1&0\\0&-1\end{pmatrix},\\
 g^0=J\otimes I_2,\quad g^1=R\otimes I_2,\quad
 g^2=Z\otimes R,\quad g^3=Z\otimes Z .
 \end{gathered}
 \label{hmtcuatro:ant:clifford}
\end{equation}
Sus cuadrados son \(-I_4,I_4,I_4,I_4\) y las matrices
distintas anticonmutan. De ello se sigue
\(\{g^I,g^J\}=2\eta^{IJ}I_4\). Definimos
\[
 g_I=\eta_{IJ}g^J,\quad
 \mathbb B_{IJ}=\tfrac14[g_I,g_J],\quad
 A_D=-ig^0,\quad \bar\psi=\psi^\dagger A_D,\quad
 \Omega=\tfrac12\omega^{IJ}\mathbb B_{IJ}.
\]
La expansión del conmutador de tres matrices da
\([\mathbb B_{IJ},g^K]=\delta_J^Kg_I-\delta_I^Kg_J\).
Por tanto \(\Omega\) representa la misma conexión métrica, con
\(D\psi=\mathrm d\psi+\Omega\psi\) y
\(D\bar\psi=\mathrm d\bar\psi-\bar\psi\Omega\).
La matriz \(A_D\) es hermítica y \(A_Dg^I\) es antihermítica;
éstas fijan el adjunto y el factor de la acción en firma \((-+++)\).

Sea \(\eta_I=\iota_{E_I}\operatorname{vol}_e\), de modo que
\(e^J\wedge\eta_I=\delta_I^J\operatorname{vol}_e\). La acción afín es
\begin{equation}
 S_D=\int_M\left\{
 \frac{\hbar}{2}
 [\bar\psi g^ID\psi-(D\bar\psi)g^I\psi]\wedge\eta_I
 -mc\,\bar\psi\psi\operatorname{vol}_e\right\}.
 \label{hmtcuatro:ant:accion-dirac}
\end{equation}
La masa y las secciones positivas \(c,\hbar\) son las lecturas
recibidas, y \([\psi]=L^{-3/2}\). La notación alternativa
\(\Gamma^I=-ig^I\), con término principal
\(i\hbar\Gamma^ID_I\), cambia a la firma de Clifford opuesta;
no se mezclan ambas convenciones.

\begin{proposition}[Corriente espinorial normalizada]
\label{hmtcuatro:ant:corriente-spin}
Con \(B^{IJK}=\bar\psi g^{[I}g^Jg^{K]}\psi\), donde la
antisimetrización está normalizada, la convención
\(\delta_\omega S_D=-\frac12\int\delta\omega^{JK}\wedge\sigma_{JK}\)
da
\begin{equation}
 \sigma_{JK}
 =-\frac{\hbar}{2}\bar\psi\{g^I,\mathbb B_{JK}\}\psi\,\eta_I
 =-\frac{\hbar}{2}B^I{}_{JK}\eta_I .
 \label{hmtcuatro:ant:sigma-spin}
\end{equation}
La corriente es independiente de la contorsión al fijar \(e,\psi\).
\end{proposition}
\begin{proof}
Se tiene
\(\delta\Omega=\frac12\delta\omega^{JK}\mathbb B_{JK}\).
Los dos términos cinéticos de \eqref{hmtcuatro:ant:accion-dirac}
dan
\[
 \delta_\omega S_D=\frac{\hbar}{4}\int
 \delta\omega^{JK}\wedge\eta_I\,
          \bar\psi\{g^I,\mathbb B_{JK}\}\psi .
\]
Comparar con la convención variacional determina signo y factor.
Al desarrollar el anticonmutador mediante Clifford, las
contracciones de grado uno se cancelan y queda
\(\{g^I,\mathbb B_{JK}\}=g^{[I}g_Jg_{K]}\).
La acción es lineal en \(\Omega\), lo que prueba la independencia.
\end{proof}

Si se incluyen conexiones internas, actúan sobre el factor
material del módulo espinorial, mientras que
\(\mathbb B_{IJ}\) actúa sobre el factor de espín. Así,
\([\,\mathbb B_{IJ}\otimes I,I\otimes\rho(A)\,]=0\).
Variar \(\omega\) produce la misma fórmula, contraída sobre
las componentes internas y sumada sobre los multipletes presentes.
Esto preserva sus representaciones quirales, de color y de familia;
no añade multipletes a los ya construidos.

\subsection{Contracción torsional y cruces entre especies}
\label{hmtcuatro:ant:contraccion}

La inversión real de Cartan--Holst y de la aplicación
\(T^I\mapsto T^I\wedge e^J-e^I\wedge T^J\), demostrada en
\ref{hmtcuatro:gea:torsion}, fija sin ambigüedad la contorsión
\(\mathfrak k_*\) para la corriente afín
\eqref{hmtcuatro:ant:sigma-spin}. En esas fórmulas se conserva
\(\gamma\in\mathbb R\setminus\{0\}\), una cotetrada no degenerada
y las secciones \(G,c,\hbar,\gamma\) fijas durante la variación.
La eliminación utiliza una sola vez la interacción lineal y da
\(-\frac14\mathfrak k_*^{IJ}\wedge\sigma_{IJ}\).

\begin{theorem}[Coeficiente de la contracción espinorial]
\label{hmtcuatro:ant:coeficiente}
Sea \(\kappa=8\pi G/c^4\). La interacción efectiva es
\[
 \mathcal L_{\rm spin}^{\rm ef}
 =C_\gamma q(B)\operatorname{vol}_e,\qquad
 C_\gamma=\frac{3\kappa c\hbar^2}{16}
                    \frac{\gamma^2}{1+\gamma^2},
\]
donde
\[
 q(B)=\frac1{3!}B^{IJK}B_{IJK}
 =-(B^{012})^2-(B^{013})^2-(B^{023})^2+(B^{123})^2 .
\]
\end{theorem}
\begin{proof}
Se calcula primero
\[
 \mathcal Y^{IJ}
 =\kappa c\frac{\gamma^2}{1+\gamma^2}
       (-\star+\gamma^{-1}I)\sigma^{IJ},\quad
 \mathcal B^I=\sum_J\iota_{E_J}\mathcal Y^{IJ},\quad
 T^I=\mathcal B^I-\tfrac14e^I\wedge\sum_L\iota_{E_L}\mathcal B^L,
\]
y después
\(\mathfrak k_{IJK}=(T_{KIJ}-T_{IJK}-T_{JKI})/2\).
Estas son composiciones lineales de las inversas ya demostradas.
Para explicitar todos los coeficientes de la contracción, se
ordena el trivector como \((012,013,023,123)\) y se extrae el factor
\(\kappa c\hbar^2\gamma^2/(1+\gamma^2)\). La sustitución de
\(\sigma_{JK}/\hbar=-\frac12B^I{}_{JK}\eta_I\) en esas composiciones
y en \(-\frac14\mathfrak k^{IJ}\wedge\sigma_{IJ}\) da
\[
 \begin{array}{c|rrrr|c}
  \text{operador sobre }\sigma&012&013&023&123&
           \text{seis coeficientes de polarización}\\ \hline
  -\star&-3/16&-3/16&-3/16&3/16&0\\
  I&0&0&0&0&0
 \end{array}
\]
Cada entrada usa \(e^I\wedge\eta_J=\delta_J^I\operatorname{vol}_e\)
y las contracciones escritas arriba. Las cuatro entradas diagonales
y las seis polarizaciones enumeran la forma cuadrática completa.
La parte \(\gamma^{-1}I\) se anula y la restante es
\((3/16)q(B)\), lo que prueba la fórmula.
\end{proof}

Para varias especies, la variación se efectúa sobre su acción
común: \(\sigma=\sum_s\sigma_s\). Como la respuesta
\(\mathfrak k_*=\mathsf C_e(\sigma)\) es lineal,
\begin{equation}
 -\tfrac14\mathsf C_e\!\left(\sum_s\sigma_s\right)
       \wedge\sum_t\sigma_t
 =-\tfrac14\sum_{s,t}\mathsf C_e(\sigma_s)\wedge\sigma_t .
 \label{hmtcuatro:ant:cruces-especies}
\end{equation}
Esta expansión demuestra que los términos \(s\ne t\) permanecen.
Calcular respuestas separadas y conservar sólo \(s=t\) cambiaría
la acción total.

\subsection{Orden normal de la corriente total}
\label{hmtcuatro:ant:car}

En la realización fermiónica finita de una celda se utiliza
\(\mathcal F=\bigoplus_n\Lambda^n V\), con las relaciones
\(\{a_r,a_s^\dagger\}=\delta_{rs}\), \(\{a_r,a_s\}=0\) y el
vacío de ocupación como referencia del orden normal.
Escribimos \(d\Gamma(M)=\sum_{r,s}a_r^\dagger M_{rs}a_s\);
esta \(\Gamma\) denota segunda cuantización, no la holonomía nonádica.
Para el módulo espinorial de cuatro componentes, las matrices son
\[
 \begin{array}{c|cccc}
 abc&012&013&023&123\\ \hline
 M_{abc}=A_Dg^ag^bg^c&
 iJ\otimes R&iJ\otimes Z&iI_2\otimes J&iZ\otimes J
 \end{array}.
\]
Multiplicar \eqref{hmtcuatro:ant:clifford} da estas matrices;
son hermíticas, de traza cero y cuadrado \(I_4\).

\begin{proposition}[Contracción normal y conservación de los cruces]
\label{hmtcuatro:ant:orden-normal}
Para toda matriz \(M\),
\[
 :d\Gamma(M)^2:
 =d\Gamma(M)^2-d\Gamma(M^2).
\]
En el módulo de cuatro componentes,
\[
 \widehat{\mathscr Q}_{\rm sp}
 =\sum_{a=1}^4s_a d\Gamma(M_a)^2+2\widehat N,
 \qquad s=(-1,-1,-1,+1),\quad \widehat N=d\Gamma(I_4).
\]
El operador es autoadjunto y conserva la ocupación.
En una suma de multipletes se emplea la misma primera identidad
con las matrices totales \(\widetilde M_a\); se conservan las
contracciones entre especies.
\end{proposition}
\begin{proof}
En \(d\Gamma(M)^2\), la relación
\(a_sa_t^\dagger=\delta_{st}-a_t^\dagger a_s\) separa la
contracción \(d\Gamma(M^2)\) del término con dos creadores y
dos aniquiladores. Este último es, por definición, el producto
normal. Como \(M_a^2=I_4\) y \(\sum_a s_a=-2\), la sustracción
de contracciones vale \(+2\widehat N\). La hermiticidad de las
matrices hace autoadjuntos los sumandos y cada monomio conserva
la ocupación.

Si \(\widetilde M=\bigoplus_s M_s\), entonces
\(d\Gamma(\widetilde M)=\sum_s d\Gamma(M_s)\) y
\(d\Gamma(\widetilde M^2)=\sum_s d\Gamma(M_s^2)\).
Los bilineales de multipletes distintos conmutan; al elevar la
primera suma al cuadrado quedan
\(2d\Gamma(M_s)d\Gamma(M_t)\) para \(s<t\). La sustracción de
una ocupación no elimina estos cruces. Es la versión operatoria
de \eqref{hmtcuatro:ant:cruces-especies}.
\end{proof}

La invariancia bajo los cambios unitarios internos se obtiene
porque éstos actúan sobre los índices materiales y entrelazan las
matrices de corriente: la segunda cuantización transporta la
fórmula completa por conjugación. La densidad y su evaluación
sobre un estado utilizan, además, el volumen de la celda y el
estado material declarados. No se identifican la corriente media
y su segundo momento. Estos datos completan los antecedentes
utilizados por la acción conjunta, sin introducir una segunda
eliminación de torsión ni modificar sus condiciones de borde.
